\subsection{\texorpdfstring{EXPANSIONS OF THE COMPLEX EXPONENTIAL, OF \(\cos x\) AND \(\sin x\).}{EXPANSIONS OF THE COMPLEX EXPONENTIAL, OF COS X AND SIN X}}

Let z be an arbitrary complex number and consider the function \(\varphi(t) = e^{zt}\) of the real variable t; we have \(\mathrm{D}^{n}\varphi(t) = z^{n}e^{zt}\) and \(e^{z} = \varphi(1)\) ; expressing \(\varphi(1)\) by means of its Taylor series of order n about the point t = 0 (II, p.~62) thus gives

\[
e ^ {z} = 1 + \frac {z}{1 !} + \frac {z ^ {2}}{2 !} + \dots + \frac {z ^ {n}}{n !} + z ^ {n + 1} \int_ {0} ^ {1} \frac {(1 - t) ^ {n}}{n !} e ^ {z t} d t\tag{6}
\]

a formula which is equivalent to (1) when \(z\) is real. The remainder

\[
r _ {n} (z) = z ^ {n + 1} \int_ {0} ^ {1} \frac {(1 - t) ^ {n}}{n !} e ^ {z t} d t
\]

in this formula can be bounded above, in absolute value, by using the inequality of the mean; if \(z = x + iy\) we have \(|e^{zt}| = e^{xt}\) , so \(|e^{zt}| \leqslant 1\) if \(x \leqslant 0\) , \(|e^{zt}| \leqslant e^x\) if \(x > 0\) ; hence

\[
\left| r _ {n} (z) \right| \leqslant \frac {\left| z \right| ^ {n + 1}}{(n + 1) !} \text {   if   } x \leqslant 0\tag{7}
\]

\[
| r _ {n} (z) | \leqslant \frac {| z | ^ {n + 1} e ^ {x}}{(n + 1) !} \text {   if   } x > 0.\tag{8}
\]

As above we conclude that

\[
e ^ {z} = \sum_ {n = 0} ^ {\infty} \frac {z ^ {n}}{n !}\tag{9}
\]

the series being absolutely and uniformly convergent on every compact subset of \(\mathbf{C}\) . From (6) one deduces in particular that

\[
e ^ {i x} = 1 + \frac {i x}{1 !} + \frac {i ^ {2} x ^ {2}}{2 !} + \dots + \frac {i ^ {n} x ^ {n}}{n !} + i ^ {n + 1} \int_ {0} ^ {x} \frac {(x - t) ^ {n}}{n !} e ^ {i t} d t\tag{10}
\]

from which we deduce the Taylor expansions of \(\cos x\) and \(\sin x\) ; on taking the real part of (10) for order \(2n + 1\) we have

\[
\cos x = 1 - \frac {x ^ {2}}{2 !} + \dots + (- 1) ^ {n} \frac {x ^ {2 n}}{(2 n) !} + (- 1) ^ {n + 1} \int_ {0} ^ {x} \frac {(x - t) ^ {2 n + 1}}{(2 n + 1) !} \cos t d t\tag{11}
\]

with remainder bounded by

\[
\left| \int_ {0} ^ {x} \frac {(x - t) ^ {2 n + 1}}{(2 n + 1) !} \cos t d t \right| \leqslant \frac {| x | ^ {2 n + 2}}{(2 n + 2) !}.\tag{12}
\]

Similarly, taking the imaginary part of (10) for order \(2n\) , we obtain

\[
\begin{array}{l} \sin x = x - \frac {x ^ {3}}{3 !} + \frac {x ^ {5}}{5 !} + \dots + (- 1) ^ {n - 1} \frac {x ^ {2 n - 1}}{(2 n - 1) !} \\ \qquad + (- 1) ^ {n} \int_ {0} ^ {x} \frac {(x - t) ^ {2 n}}{(2 n) !} \cos t d t \end{array}\tag{13}
\]

with remainder bounded by

\[
\left| \int_ {0} ^ {x} \frac {(x - t) ^ {2 n}}{(2 n) !} \cos t d t \right| \leqslant \frac {| x | ^ {2 n + 1}}{(2 n + 1) !}.\tag{14}
\]

Moreover, on comparing the remainders in (11) for orders \(2n + 1\) and \(2n + 3\) , we have

\[
\int_ {0} ^ {x} \frac {(x - t) ^ {2 n + 3}}{(2 n + 3) !} \cos t d t = \frac {x ^ {2 n + 2}}{(2 n + 2) !} - \int_ {0} ^ {x} \frac {(x - t) ^ {2 n + 1}}{(2 n + 1) !} \cos t d t
\]

and taking (12) into account we see that the remainder in (11) has the same sign as \((-1)^{n+1}\) no matter what x; in the same way we can show that the remainder in (13) has the same sign as \((-1)^{n}x\) . In particular, for n = 0 and n = 1 in (11), and for n = 1 and n = 2 in (13) we obtain the inequalities

\[
1 - \frac {x ^ {2}}{2} \leqslant \cos x \leqslant 1 \quad \text { for   all } x\tag{15}
\]

\[
x - \frac {x ^ {3}}{6} \leqslant \sin x \leqslant x \quad \text { for   all } x \geqslant 0.\tag{16}
\]

Finally, on putting z = ix in (9) we have

\[
\cos x = \sum_ {n = 0} ^ {\infty} (- 1) ^ {n} \frac {x ^ {2 n}}{(2 n) !}\tag{17}
\]

\[
\sin x = \sum_ {n = 0} ^ {\infty} (- 1) ^ {n} \frac {x ^ {2 n + 1}}{(2 n + 1) !}\tag{18}
\]

these series being absolutely and uniformly convergent on every compact interval.

It is clear, furthermore, that the formulae (17) and (18) remain valid for every complex x, the series on the right-hand side being absolutely and uniformly convergent on every compact subset of C. In particular, for every x (real or complex)

\[
\begin{array}{l} \cosh x = \sum_ {n = 0} ^ {\infty} \frac {x ^ {2 n}}{(2 n) !} \\ \sinh x = \sum_ {n = 0} ^ {\infty} \frac {x ^ {2 n + 1}}{(2 n + 1) !}. \end{array}
\]

